\scanpage{15}
\mspage{13}
comme il est légitime, trouver un analogue \deleted{pour} \inserted{\unsure{de}} la restriction à $X_K$ \inserted{\unsure{envisagée|envisagé}} plus haut, il faut définir $X_K$ comme \uline{l'espace rigide-analytique sur $K$ défini par le schéma formel} $X$, et considérer sur $X_K$ des Modules cohérents munis de \deleted{connexions} \uline{stratifications au sens rigide-analytique}, ou ce qui revient au même, de connexions intégrables \inserted{en ce sens}. De tels Modules sont encore nécessairement localement libres. On trouve ainsi un foncteur des cristaux de Modules cohérents sur $S$ dans les Modules cohérents stratifiés sur $X_K$, induisant un foncteur pleinement fidèle de la catégorie des isocristaux cohérents sur $S$, dans la catégorie des Modules cohérents stratifiés sur $X_K$. Comme tout à l'heure (et même plus, si on peut dire, car c'est vraiment \deleted{\unsure{si|ici}} ici la situation ``naturelle''), il s'impose de regarder quelle est l'image essentielle. On peut également se demander si les Modules cohérents stratifiés sur un espace rigide-analytique n'admettraient pas une description simple, en termes d'un groupe plus ou moins discret jouant le rôle du groupe fondamental dans la théorie transcendante sur le corps des complexes.

\deleted{\unsure{1.14.}} La description donnée \inserted{des cristaux sur $S$}, toute triviale qu'elle soit, a déjà des conséquences intéressantes pour la structure de la catégorie des cristaux \inserted{de modules} cohérents sur $S$: \deleted{\unsure{un tel cristal}} cette catégorie est noethérienne (si $S$ est noethérien i.e. de type fini sur $k$), tout objet contient donc un plus grand sous-objet de torsion, de plus \deleted{\unsure{un}} \inserted{les} objets sans torsion correspond\inserted{ent} dans la description ci-dessus \deleted{\unsure{à un}} \inserted{aux} Modules stratifiés sans torsion sur $X$, lesquels \deleted{\unsure{est}} \inserted{sont} alors automatiquement localement libres. Il est bien probable que des résultats analogues doivent être vrais, sans hypothèse du genre lissité sur $S$.

1.14. Il faut encore expliciter, dans la description générale précédente, la \deleted{donnée} foncteur $M\rightsquigarrow M^{(p)}$, pour pouvoir décrire de façon analogue les bicristaux et bi-isocristaux. De façon générale, si en plus de la donnée de $S$ sur $k$, on a un $S'$ lisse sur $k'$ parfait, et des \deleted{\unsure{homo}}morphismes $S'\longrightarrow S$ et $\operatorname{Spec}(k')\longrightarrow\operatorname{Spec}(k)$ donnant lieu à un carré commutatif, et si enfin on peut relever $S'$ formellement en $X'$, et $S'\longrightarrow S$ en $X'\longrightarrow X$, donnant un carré commutatif avec $\operatorname{Spec}W'\longrightarrow\operatorname{Spec}W$, alors la notion d'image inverse de cristaux de modules cohérents, de $S$ à $S'$, s'explicite en termes d'image inverse de Modules cohérents stratifiés, de $X$
\margin{14}
